\documentclass[11pt,a4paper]{article}

\usepackage[utf8]{inputenc}
\usepackage[T1]{fontenc}
\usepackage[ngerman]{babel}
\usepackage{lmodern}
\usepackage{amsmath,amssymb}
\usepackage{graphicx}
\usepackage{booktabs,tabularx,array}
\usepackage[table]{xcolor}
\usepackage{tikz}
\usetikzlibrary{arrows.meta,positioning,shapes.geometric,calc}
\usepackage[most]{tcolorbox}
\usepackage{enumitem}
\usepackage{microtype}
\usepackage{hyperref}
\usepackage[a4paper,left=18mm,right=18mm,top=15mm,bottom=15mm,headheight=0pt]{geometry}

\definecolor{navy}{HTML}{174A63}
\definecolor{blue}{HTML}{2E7698}
\definecolor{lightblue}{HTML}{EAF3F7}
\definecolor{lightgray}{HTML}{F2F3F4}
\definecolor{midgray}{HTML}{66727A}
\definecolor{green}{HTML}{477A63}
\definecolor{orange}{HTML}{B96A2B}

\hypersetup{colorlinks=true,linkcolor=navy,urlcolor=blue,pdftitle={Sicherungsblatt Aufgabe 3b -- Klassifikation der Testdaten}}
\setlength{\parindent}{0pt}
\setlength{\parskip}{4pt}
\setlist[itemize]{leftmargin=5mm,itemsep=1.5pt,topsep=2pt}
\setlist[enumerate]{leftmargin=6mm,itemsep=2pt,topsep=2pt}
\renewcommand{\familydefault}{\sfdefault}

\pagestyle{empty}

\newtcolorbox{definition}[1][]{enhanced,breakable,colback=lightblue,colframe=blue,
  boxrule=0.8pt,arc=1.5mm,left=3mm,right=3mm,top=2mm,bottom=2mm,
  fonttitle=\bfseries,title=Definition,#1}
\newtcolorbox{merke}[1][]{enhanced,breakable,colback=lightgray,colframe=navy,
  boxrule=0.9pt,arc=1.5mm,left=3mm,right=3mm,top=2mm,bottom=2mm,
  fonttitle=\bfseries,title=Merke,#1}
\newtcolorbox{beispiel}[1][]{enhanced,breakable,colback=white,colframe=midgray,
  boxrule=0.7pt,arc=1.5mm,left=3mm,right=3mm,top=2mm,bottom=2mm,
  fonttitle=\bfseries,title=Beispiel,#1}

\newcommand{\fach}[1]{\textbf{\color{navy}#1}}
\newcommand{\sectionline}[1]{%
  \vspace{1mm}{\Large\bfseries\color{navy}#1}\par\vspace{1mm}%
  {\color{blue}\rule{\linewidth}{0.8pt}}\vspace{1mm}}
\newcommand{\sheettitle}[2]{%
  \begin{tcolorbox}[colback=navy,colframe=navy,arc=2mm,boxrule=0pt,left=5mm,right=5mm,top=2.5mm,bottom=2.5mm]
    {\color{white}\fontsize{20}{24}\selectfont\bfseries #1}\par
    \vspace{0.5mm}{\color{white}\large #2}
  \end{tcolorbox}}

\definecolor{fishblue}{HTML}{377ED1}
\definecolor{fishorange}{HTML}{E87931}
\definecolor{fishyellow}{HTML}{F2C94C}
\definecolor{fishred}{HTML}{D84B53}
\definecolor{fishdark}{HTML}{23324A}
\definecolor{peaceful}{HTML}{08705A}
\definecolor{hostile}{HTML}{992D35}

% Wie in Sicherungsblatt 3a; das Label unter dem Fisch ist das tatsaechliche Label.
% \fisch{(Position)}{Schuppenfarbe}{Punkte 0/1}{Bauchfarbe}{Flossenfarbe}{ID}{Label f/h}{Fehler 0/1}
\newcommand{\fisch}[8]{%
  \begin{scope}[shift={#1}]
    \filldraw[fill=#5,draw=fishdark,line width=0.5pt,line join=round] (0.46,0) -- (0.74,0.21) -- (0.74,-0.21) -- cycle;
    \filldraw[fill=#5,draw=fishdark,line width=0.5pt,line join=round] (-0.12,0.21) -- (0.08,0.36) -- (0.2,0.19) -- cycle;
    \begin{scope}
      \clip (0,0) ellipse (0.5 and 0.26);
      \fill[#2] (-0.6,-0.3) rectangle (0.6,0.3);
      \fill[#4] (-0.6,-0.3) rectangle (0.6,-0.1);
      \ifnum#3=1
        \fill[fishdark] (-0.12,0.09) circle (0.035) (0.1,0.14) circle (0.035) (0.28,0.05) circle (0.035)
          (0.05,-0.01) circle (0.035) (-0.2,-0.03) circle (0.035);
      \fi
    \end{scope}
    \draw[fishdark,line width=0.7pt] (0,0) ellipse (0.5 and 0.26);
    \filldraw[fill=fishdark,draw=white,line width=0.4pt] (-0.32,0.06) circle (0.055);
    \if#7f
      \node[font=\tiny,text=peaceful,align=center,inner sep=0pt] at (0.12,-0.66) {\textbf{#6}\\friedlich};
    \else
      \node[font=\tiny,text=hostile,align=center,inner sep=0pt] at (0.12,-0.66) {\textbf{#6}\\feindselig};
    \fi
    \ifnum#8=1
      \draw[fishred,line width=0.9pt,dashed,rounded corners=1mm] (-0.6,-0.98) rectangle (0.84,0.5);
      \node[fill=fishred,text=white,font=\tiny\bfseries,rounded corners=0.8mm,inner sep=1pt] at (0.12,0.5) {Fehler};
    \fi
  \end{scope}}

\newcommand{\zelle}[3]{\textbf{\large #1}\quad{\footnotesize #2}\newline{\scriptsize\color{midgray}#3}}

\begin{document}

\sheettitle{Klassifikation der Testdaten}{Entscheidungsbäume · Sicherungsblatt zu Aufgabe 3b · Informatik 11}

\sectionline{1. Testdaten klassifizieren}

\fach{Testdaten} wurden nicht zum Erstellen des Baums verwendet. Ihr tatsächliches Label ist bekannt.

\begin{center}
\begin{tikzpicture}[font=\small,every node/.style={align=center},
  decision/.style={draw=navy,fill=lightblue,rounded corners=1mm,minimum width=27mm,minimum height=8mm},
  leaf/.style={draw=midgray,fill=lightgray,rounded corners=1mm,minimum width=21mm,minimum height=8mm},
  group/.style={draw=midgray,fill=white,rounded corners=2mm,line width=0.7pt}]
  \node[decision] (root) at (0,0) {Schuppenfarbe?};
  \node[decision] (pattern) at (-4.2,-1.4) {Muster?};
  \node[decision] (bellyorange) at (4.6,-1.4) {Bauchfarbe?};
  \node[leaf] (none) at (-6.6,-2.8) {friedlich};
  \node[decision] (bellyblue) at (-1.8,-2.8) {Bauchfarbe?};
  \node[leaf] (bbBlack) at (-3.4,-4.2) {feindselig};
  \node[leaf] (bbWhite) at (-0.2,-4.2) {friedlich};
  \node[leaf] (obBlack) at (3.0,-2.8) {friedlich};
  \node[leaf] (obWhite) at (6.2,-2.8) {feindselig};
  \draw (root) -- node[above left,font=\footnotesize]{Blau} (pattern);
  \draw (root) -- node[above right,font=\footnotesize]{Orange} (bellyorange);
  \draw (pattern) -- node[above left,font=\footnotesize]{Ohne} (none);
  \draw (pattern) -- node[above right,font=\footnotesize]{Punkte} (bellyblue);
  \draw (bellyblue) -- node[above left,font=\footnotesize]{Schwarz} (bbBlack);
  \draw (bellyblue) -- node[above right,font=\footnotesize]{Weiß} (bbWhite);
  \draw (bellyorange) -- node[above left,font=\footnotesize]{Schwarz} (obBlack);
  \draw (bellyorange) -- node[above right,font=\footnotesize]{Weiß} (obWhite);

  % Testfische unter dem Blatt, in dem sie ankommen
  \draw[group] (-8.05,-3.35) rectangle (-5.15,-5.2);
  \draw (none.south) -- (-6.6,-3.35);
  \fisch{(-7.45,-4.05)}{fishblue}{0}{fishdark}{fishred}{T1}{f}{0}
  \fisch{(-6.0,-4.05)}{fishblue}{0}{white}{fishred}{T5}{f}{0}

  \draw[group] (1.55,-3.35) rectangle (4.45,-5.2);
  \draw (obBlack.south) -- (3.0,-3.35);
  \fisch{(2.15,-4.05)}{fishorange}{0}{fishdark}{fishyellow}{T3}{f}{0}
  \fisch{(3.6,-4.05)}{fishorange}{1}{fishdark}{fishyellow}{T4}{h}{1}

  \draw[group] (5.4,-3.35) rectangle (7.0,-5.2);
  \draw (obWhite.south) -- (6.2,-3.35);
  \fisch{(6.08,-4.05)}{fishorange}{1}{white}{fishred}{T2}{h}{0}
\end{tikzpicture}

{\scriptsize\color{midgray}Blatt = \textbf{Vorhersage}, darunter die dort ankommenden Testfische mit \textbf{tatsächlichem Label}.
Körper = Schuppenfarbe, Punkte = Muster, Streifen = Bauchfarbe, Flossen = Flossenfarbe.}
\end{center}

\sectionline{2. Konfusionsmatrix}

\begin{minipage}[t]{0.63\linewidth}
\vspace{0pt}
\renewcommand{\arraystretch}{1.15}\setlength{\tabcolsep}{4pt}
\begin{tabularx}{\linewidth}{@{}>{\raggedright\arraybackslash}p{19mm}XX>{\centering\arraybackslash}p{13mm}@{}}
\toprule
{\scriptsize\textbf{erwartet} $\backslash$\newline\textbf{berechnet}} & \textbf{friedlich} & \textbf{feindselig} & \textbf{Summe}\\
\midrule
\textbf{friedlich}  & \cellcolor{peaceful!12}\zelle{3}{richtig positiv}{T1, T3, T5} & \cellcolor{fishred!10}\zelle{0}{falsch negativ}{--} & 3\\
\textbf{feindselig} & \cellcolor{fishred!10}\zelle{1}{falsch positiv}{T4} & \cellcolor{peaceful!12}\zelle{1}{richtig negativ}{T2} & 2\\
\midrule
\textbf{Summe} & \multicolumn{1}{c}{4} & \multicolumn{1}{c}{1} & \textbf{5}\\
\bottomrule
\end{tabularx}
\end{minipage}\hfill
\begin{minipage}[t]{0.33\linewidth}
\vspace{1mm}
\small\raggedright
\textbf{Zeile:} erwartetes, also tatsächliches Label

\textbf{Spalte:} berechnetes Label, also die Vorhersage

\medskip
\fach{positiv} = friedlich\\
\fach{negativ} = feindselig

\medskip
{\color{peaceful}\rule[-0.5mm]{3mm}{3mm}}\ Diagonale: richtig klassifiziert
\end{minipage}

\sectionline{3. Genauigkeit}

\begin{definition}[title=Definition: Genauigkeit]
Die \fach{Genauigkeit} ist der Anteil richtig klassifizierter Testdaten an allen Testdaten.
\[
  \text{Genauigkeit}=\frac{\text{richtig positiv}+\text{richtig negativ}}
  {\text{richtig positiv}+\text{falsch negativ}+\text{falsch positiv}+\text{richtig negativ}}
\]
\begin{center}
\begin{tikzpicture}[font=\scriptsize]
  \foreach \id/\x in {T1/0,T2/1,T3/2,T5/3}
    \node[draw=peaceful,fill=peaceful!15,minimum width=12mm,minimum height=6mm] at (1.25*\x,0) {\id\ \checkmark};
  \node[draw=fishred,fill=fishred!12,minimum width=12mm,minimum height=6mm] at (5,0) {T4 \texttimes};
  \node[anchor=west,font=\small] at (5.8,0) {$\displaystyle\Rightarrow\ \frac{3+1}{3+0+1+1}=\frac{4}{5}=\mathbf{80\,\%}$};
\end{tikzpicture}
\end{center}
\end{definition}

\begin{merke}
Die Qualität eines Entscheidungsbaums wird mit Testdaten geprüft, nicht mit den Trainingsdaten.\par
Die Genauigkeit beschreibt nur die verwendeten Testdaten. Sie ist keine Garantie für alle künftigen Daten.
\end{merke}

\end{document}
